\section{Servicing decisions: intervention histories and net value}
\label{exec14:l5-servicing}

The policy functional in equation~\eqref{eq:la04-policy} becomes operationally
useful only when its histories have a banking interpretation. Consider an
initial servicing contact $A_0$, a subsequent distress indicator $L_1$, and a
second contact $A_1$. An initial contact can prevent distress; distress can
prompt a second contact; and both contacts incur costs. A request to compare
``contacted and uncontacted customers'' conceals all three relations. It also
conceals whether the analyst wants the value of a fixed history, a rule that
responds to distress, or a policy learned from the same observations. These are
different inferential requests. An assistant can make the alternatives explicit
and translate the selected request into a rule. Its interpretation remains
reviewable against the business decision.

The following mechanism is stipulated for an independent exact reference; its
numbers are illustrative USD per account, not observations from a bank
\cite{exec14_workflow_oracle}. There is no baseline covariate in this small law.
The first action has probability $1/2$, and the state transitions satisfy
\[
 P(L_1=1\mid A_0=0)=\frac45,\qquad
 P(L_1=1\mid A_0=1)=\frac25.
\]
The observational second-action probabilities are
\[
 e_1(l)=P(A_1=1\mid A_0,L_1=l)
 =\begin{cases}1/2,&l=0,\\3/4,&l=1.\end{cases}
\]
Gross recovery is $100-40L_1+10A_0+30L_1A_1$. Each contact costs eight, so the
outcome to be compared is
\[
 Y=100-40L_1+2A_0+(30L_1-8)A_1.
 \label{eq:exec14-l5-net}
\]
The stipulated mechanism has independent assignment noise at each action,
conditional on the displayed history, and a deterministic final outcome.
This supplies the exchangeability needed for its illustrative intervention
interpretation. An observed bank table with the same conditional means would
still require an assignment argument; numerical resemblance would not supply
that premise. Treatment-confounder feedback and the causal conditions of the
g-formula are developed by Hern\'an and Robins \cite{hernan2020}.

Compare three fixed, feasible policies. The no-contact policy sets both
actions to zero. The treat-all policy sets both to one. The risk-responsive
policy sets $A_0=1$ and $A_1=L_1$. Its second contact is affordable and available
only after the distress state is delivered. The policy does not use a future
outcome or a revised state unavailable at that decision. Under the risk policy,
net recovery is 102 when $L_1=0$ and 84 when $L_1=1$. Averaging under its own
post-first-action state law gives
\[
 V(g_{\mathrm{risk}})=\frac35(102)+\frac25(84)
 =\frac{474}{5}=94.8.
\]
The corresponding values are $V(g_{\mathrm{none}})=68$ and
$V(g_{\mathrm{all}})=90$. The risk-responsive policy gains
$134/5=26.8$ over no contact and 4.8 over contacting everyone. The first contrast
decomposes into two from the first contact's direct net contribution, sixteen
from the reduction in distress probability, and $44/5=8.8$ from the second
contact's net contribution under the changed state distribution. This
decomposition explains why a recovery lift alone is insufficient for the
business decision. A contact can improve gross recovery while its marginal
net contribution is negative in a low-risk state.

Two plausible shortcuts answer different questions. Restricting the observed
data to histories that happen to follow the risk policy gives mean 93, because
the two retained states have equal conditional frequency. Restricting to
histories with both actions zero gives $220/3$, because distress changes the
chance of avoiding the second contact. The unweighted difference is therefore
$59/3$, rather than $134/5$. Alternatively, averaging the risk-policy outcome
under the pooled distress frequency $3/5$ yields 91.2. That calculation holds
the mediator distribution at its observational mixture, rather than allowing
the first intervention to change it. Neither error is repaired by producing
an explanation of a correctly computed mean. The assistant's most valuable
intervention is to detect the target substitution before numerical execution.

For a general baseline $L_0$, write
$H_1=(L_0,A_0,L_1)$ and
$m(H_1,a_1)=E[Y\mid H_1,A_1=a_1]$. With a fixed regime $g$, define
\begin{align*}
 m_g(H_1)&=m\{H_1,g_1(H_1)\},\\
 Q_0(L_0)&=E[m_g(H_1)\mid L_0,A_0=g_0(L_0)].
\end{align*}
Then $V(g)=E[Q_0(L_0)]$ under the stated longitudinal identification
conditions. Backward outcome regression implements these successive
conditional expectations. Forward g-computation models the state transition
and outcome, then propagates the regime through both. Inverse-probability
weighting instead uses
\[
 w_0=\frac{\mathbf1\{A_0=g_0(L_0)\}}
 {P\{A_0=g_0(L_0)\mid L_0\}},\qquad
 w_1=\frac{\mathbf1\{A_1=g_1(H_1)\}}
 {P\{A_1=g_1(H_1)\mid H_1\}},
\]
for which $E[w_0w_1Y]=V(g)$. All three representations use the intervention's
history. They differ in which nuisance quantities must be estimated and in
their behavior when data are sparse, weights are variable, or a regression
extrapolates. The choice should follow the declared target and its supported
histories, rather than the name of a preferred package.

\section{Retention, recursive augmentation and uncertainty}
\label{exec14:l5-retention}

The outcome may become unavailable when an account moves to another servicing
system or its case is unresolved. Let $R$ indicate that $Y$ is observed at the
final horizon. To identify the same full-population value, suppose
$R\perp Y\mid H_1,A_1$, with
$r(H_1,A_1)=P(R=1\mid H_1,A_1)>0$ on the regime's support. This is a final-outcome
missingness law. If earlier histories also disappear, their observation
mechanisms require additional nodes and weights; the single $R$ does not
silently represent them. Nor does administrative censoring justify this
conditional independence merely because its date is recorded
\cite{hernan2020}.

In the stipulated servicing law, let retention be one for $L_1=0$ and one half
for $L_1=1$, independently of the outcome given history. Among retained accounts
under the risk intervention, the mean becomes
\[
 \frac{(3/5)102+(2/5)(1/2)84}{3/5+(2/5)(1/2)}=97.5.
\]
Inverse retention weights restore the intervention value 94.8. Multiplying
$w_0w_1$ by $R/r$ also restores it in the observational enumeration. The
complete-case mean is higher because distressed accounts are less often
retained, not because servicing became more effective. A fluent report that
describes 97.5 as a population policy value would obscure an observation-law
error. Conversely, if loss depends on an unobserved recovery component even
after conditioning, these weights have no such causal justification. The
appropriate output is a sensitivity analysis or unresolved population value,
not a fabricated retention probability.

An augmented representation shows where the pieces enter. With the true
nuisance quantities, a centered score for this two-action, final-retention
problem is
\begin{align}
 \phi_g(O)={}&w_0w_1\frac{R}{r(H_1,A_1)}
          \{Y-m(H_1,A_1)\}\nonumber\\
 &+w_0\{m_g(H_1)-Q_0(L_0)\}+Q_0(L_0)-V(g).
 \label{eq:exec14-l5-score}
\end{align}
The first residual has conditional mean zero under the observation law. The
second has mean zero conditional on $L_0$ after first-action reweighting.
The last term centers the baseline regression. Where $R=0$, the first term
is defined as zero; it does not require an unavailable value of $Y$. This
mean-zero argument is a population identity. A fitted estimator needs control
of nuisance errors and remainder terms, a valid sampling law, and the
appropriate handling of sample splitting. The display does not establish
coverage simply by being called an influence function.

The exact law isolates one aspect of precision. Its outcome is deterministic
given history and it has no varying $L_0$. The first and last components in
equation~\eqref{eq:exec14-l5-score} vanish, leaving variation from the intermediate
state. Under independent accounts, the risk-policy score has variance
$3888/25=155.52$; the no-contact score has variance 512. Their covariance is
zero here because their nonzero first-action indicators are disjoint and both
scores are centered. The variance of their difference is consequently
$16688/25=667.52$. A normalized policy-IPW estimator with true weights has
asymptotic centered score $w_0w_1\{Y-V(g)\}$ when there is no retention loss.
Its variance is 248.832 for the risk policy and 1228.8 for no contact.
These are exact moments under the declared law
\cite{exec14_workflow_oracle}. They explain how recursive adjustment can remove
second-action sampling variation in this example. They do not rank fitted
procedures under arbitrary bank dependence or qualify any confidence interval.

Robustness requires equally precise terminology. Luedtke and colleagues
distinguish conventional longitudinal double robustness, requiring all outcome
components or all treatment components to be correct in the relevant sense,
from sequential double robustness, allowing the appropriate component to be
correct at each stage \cite{sequentialdr2018}. Their definitions distinguish
outcome-regression specification from treatment-mechanism consistency. A
point-treatment AIPW function placed inside a loop does not inherit their
recursive construction or its theorem. An assistant recommending a
``doubly robust longitudinal estimator'' should therefore identify the actual
algorithm and its conditions, rather than treating robustness as a universal
package property. The safe comparison includes a transparent fixed-regime
g-formula, an IP-weighted representation, and an estimator whose longitudinal
construction is genuinely supported by primary text.

Support is a property of the policy's histories, not only a histogram of the
first propensity. An initially supported contact can lead to a distressed
state where the requested later action is never taken. Product probabilities
can also become very small although each decision separately looks tolerable.
Weight truncation may improve numerical behavior while changing bias and the
effective comparison. Diagnostics should show where weights concentrate,
which histories disappear, and how the result changes under a declared
sensitivity choice. A regime that requires a forbidden or unavailable contact
should be returned for redesign. A new policy learned after inspecting the
same outcomes introduces selection and policy-learning questions beyond the
fixed-policy arithmetic above. Its apparent best value cannot be reported
with a fixed-regime interval without a separate argument.

\section{Branch adoption is a panel question with a different target}
\label{exec14:l5-panel}

Suppose branches adopt a servicing platform at different dates and the bank
asks whether the platform improves recovery. This differs from evaluating the
account-level dynamic contact rule. Let $G_i$ be branch $i$'s first adoption
time, $Y_{it}(g)$ its outcome at time $t$ if adoption begins at $g$, and
$Y_{it}(\infty)$ the outcome without adoption. A useful target is
\[
 \operatorname{ATT}(g,t)
 =E[Y_{it}(g)-Y_{it}(\infty)\mid G_i=g],\qquad t\geq g.
\]
The unit, adoption event, composition of branch accounts, and monetary outcome
must be fixed before estimation. A decline in average loss among remaining
accounts could arise from account transfers rather than an improvement for
the original cohort. A repeated cross-section and a stable account panel
therefore need different observation arguments. Arkhangelsky and Imbens
organize panel methods by data configuration, relative numbers of units and
periods, and treatment assignment patterns; this organization helps avoid
selecting a method solely because the table has a date column
\cite{panelreview2024}.

With no anticipation and an appropriate never-treated comparison group $C$,
unconditional parallel untreated trends over the relevant periods identify
\begin{align*}
 \operatorname{ATT}(g,t)
 ={}&E[Y_{it}-Y_{i,g-1}\mid G_i=g]\\
 &-E[Y_{it}-Y_{i,g-1}\mid C_i=1].
\end{align*}
Conditional parallel trends instead require comparing changes at the relevant
baseline covariates $X$ and averaging over the adopting cohort's covariate
distribution, with overlap for that comparison. Callaway and Sant'Anna also
develop not-yet-treated comparisons and explicit anticipation windows
\cite{callaway2021}. For an anticipation horizon $\delta$, the baseline must
precede $g-\delta$, and a comparison group must still be unaffected at the
relevant later time including its anticipation window. A branch preparing for
rollout can change staffing before formal adoption; the recorded software
installation date alone does not establish no anticipation.

The comparison choice carries substantive assumptions. Never-treated branches
may operate in different markets. Not-yet-treated branches may be more similar
but supply no untreated comparison after the last adoption, and their rollout
plans may already affect behavior. Both designs require an argument about
untreated changes, rather than the stronger sequential exchangeability of an
account-level dynamic policy. The analyst must state which argument is being
used. An assistant can locate adoption notices, reconstruct candidate dates,
and propose competing controls. It cannot establish the unobserved untreated
trend from a satisfactory regression fit.

A small stipulated panel makes these distinctions visible
\cite{exec14_workflow_oracle}. There are three equal-sized cohorts, early
adopters $E$, late adopters $L$, and never adopters $N$. Their observed means are
\begin{center}
\begin{tabular}{lrrr}
\toprule
Cohort & Time 0 & Time 1 & Time 2\\
\midrule
Early &100&107&119\\
Late &100&102&108\\
Never &100&102&104\\
\bottomrule
\end{tabular}
\end{center}
The stipulated untreated means are 100, 102, and 104 for every cohort.
Early adoption has effects five at time 1 and fifteen at time 2; late adoption
has effect four at time 2. Comparing early and never cohorts from time 0 to
time 1 gives five; from time 0 to time 2 gives fifteen. Comparing late and
never cohorts from time 1 to time 2 gives four. In contrast, comparing the
late cohort's change to the already-treated early cohort's change gives
\[
 (108-102)-(119-107)=-6.
\]
The early cohort's treatment effect increases by ten during that comparison.
The negative difference therefore subtracts another group's changing effect
from the late cohort's benefit. It is not evidence that the platform harms
late adopters.

Aggregation also changes the question. Equal weighting of the three treated
cohort-time cells gives $(5+15+4)/3=8$. Equal weighting of the two cohorts at
their first exposure gives $(5+4)/2=4.5$. The first includes a later-exposure
effect for early adopters; the second compares a common event horizon. Other
bank questions may weight cohort sizes, account exposure, or the distribution
of dates in a proposed rollout. These weights must be stated. In this particular
balanced fixture, residualizing the adoption indicator against cohort and
calendar effects produces a two-way fixed-effects coefficient of 4.5: the
early time-2 cell has zero residual weight and the other treated cells have
equal positive weights. It happens to equal the first-exposure target here.
The contaminated two-group comparison does not imply that every fixed-effects
coefficient is negative or scientifically useless. The exact residualization
and the survey's qualified discussion both argue for inspecting the target and
weights, rather than adopting a universal rejection of fixed effects
\cite{panelreview2024,exec14_workflow_oracle}.

Observational ambiguity remains even after this algebra is correct. Suppose
the early cohort's untreated outcome at time 2 is 107 because it receives a
cohort-specific favorable shock of three, and its actual treatment effect
is twelve. The observed value remains 119, and every other displayed value is
unchanged. The first law's effect fifteen and the alternative law's effect
twelve are compatible with the same observed panel. The alternative violates
the relevant post-adoption parallel untreated trend. Additional common
pre-adoption periods can be appended to both laws without distinguishing
them. Good pre-adoption agreement is a useful diagnostic; it does not exclude
every later counterfactual shock. A bank may need a rollout assignment design,
a defensible control, or explicit sensitivity to such shocks before attaching
a causal interpretation to the contrast.

Synthetic control and factor-based alternatives are substantive when the
bank has few treated units, rich histories, and a credible structure for
untreated outcomes. They do not automatically solve an unsupported
counterfactual: donor contamination, a new factor after rollout, or a change
in measurement can defeat their argument. The panel survey motivates these
alternatives without asserting a particular bank guarantee
\cite{panelreview2024}. The assistant's proposal should connect each candidate
to the data dimensions and assignment pattern, identify which restriction
supplies its counterfactual, and retain uncertainty about that restriction.
A dynamic servicing value should not be silently substituted for a branch
adoption ATT merely because both methods return a monetary number.

\section{Observation construction can determine the answer}
\label{exec14:l5-observations}

The preceding values assume that accounts and outcomes have been correctly
assembled. A causal workflow should instead begin with the population frame,
inclusion mechanism, ownership relation, and label meaning. An assistant
extracting ``recovery'' from notes may confuse gross recovery with net value,
or a revised figure with the figure available at the decision. Linking two
accounts to a customer may change the unit, create duplicated outcomes, and
alter dependence at once. These changes belong to the scientific observation
model. A schema that accepts the resulting rows does not establish their truth.

Consider a frame containing two equally frequent strata with values zero and
100. Suppose every zero-value unit is included but each 100-value unit is
included independently with probability $1/4$. Writing $S$ for inclusion,
the mean among observed records is
\[
 \frac{E[SY]}{E[S]}=
 \frac{(1/2)(1/4)100}{(1/2)1+(1/2)(1/4)}=20.
\]
The population mean is 50. If the frame size $N$ and truthful inclusion
probabilities $\pi_i$ are known, the Horvitz--Thompson form
\[
 \widehat\mu_{\mathrm{HT}}=
 \frac1N\sum_{i=1}^N\frac{S_iY_i}{\pi_i}
\]
has expectation 50 under the stipulated sampling law. This uses the known
frame denominator; a ratio estimator with an estimated denominator is a
different construction. For independently generated framed units, the
per-unit weighted contribution is zero or 400. Its variance is
$20000-50^2=17500$, so the mean's variance is $17500/N$
\cite{exec14_workflow_oracle}. Correcting selection can therefore increase
uncertainty even while restoring the target. A narrow unweighted interval
around 20 does not constitute a more precise estimate of 50.

The correction requires knowledge that an LLM cannot manufacture. If inclusion
probabilities are estimated from an incomplete frame, depend on an unknown
outcome component, or are zero for a relevant group, the displayed equality
need not apply. Retrieving a description of the sampling system can identify
a candidate probability model; an actual audit must establish that it was
followed. When the target is an intervention mean, treatment, retention,
sampling, and linkage cannot simply be handled as unrelated cleaning steps.
Their joint chronology determines which histories and outcomes are observed
and which probabilities can be conditioned on. The longitudinal observation
intervention discussed by Hern\'an and Robins supplies a model-based way to
formulate that joint problem \cite{hernan2020}.

Noisy labels pose a separate problem. Let $T$ be a true binary ownership or
distress label, $Z$ its observed classifier label, $p=P(T=1)$, sensitivity
$s=P(Z=1\mid T=1)$ and specificity $c=P(Z=0\mid T=0)$. If these error rates are
known for the target population, then
\[
 P(Z=1)=sp+(1-c)(1-p),\qquad
 p=\frac{P(Z=1)+c-1}{s+c-1}.
\]
With $p=0.4$, $s=0.8$ and $c=0.9$, the positive-label frequency is 0.38 and
the inversion returns 0.4. If $s+c=1$, the observed label frequency contains
no information about prevalence through this model. Close to that boundary,
sampling noise is amplified by the denominator. With known rates and $n$
independent labels in this example, the inverted proportion has variance
$(589/1225)/n$ before any probability-bound handling
\cite{exec14_workflow_oracle}. Estimating the error rates from a gold-label
sample adds uncertainty. Reusing one error-rate estimate across many account
records also introduces shared estimation error; multiplying the number of
rows does not make that component disappear.

This inversion concerns a marginal prevalence, not a causal effect or an
individual ownership map. Errors can differ by treatment, history, market or
time, and two competing customer mappings can have the same correctness
count. A prevalence correction does not reconstruct which record belongs to
which economic owner. For causal use, the workflow must retain the candidate
mapping and how its uncertainty affects exposure, outcome and unit formation.
An LLM can propose label rules and flag contradictory notes; the proposal
should compare that assistance with deterministic rules and human adjudication
on genuinely reviewed examples. Agreement between two assistants exposed to
the same notes is not an independent gold label. Missing truth is an evidence
limitation, rather than a reason to replace adjudication with confidence in
generated prose.

Dependence can survive correct inclusion and label handling. Let $B$ independent
branches each have a centered shock $\varepsilon_b$ of variance one, shared
by two account records. Their overall account mean is
$B^{-1}\sum_b\varepsilon_b$, with variance $1/B$. Treating the $2B$ rows as
independent instead gives $1/(2B)$. Now add a centered calendar shock $C$ of
variance four shared by every branch, independently of their shocks. The
mean's variance is $4+1/B$. More accounts or branches reduce only the second
component. This is an exact covariance construction
\cite{exec14_workflow_oracle}, not a universal banking variance model. It
shows why account counts alone cannot justify precision and why a common
calendar component needs replicated dates or additional restrictions.
Splitting nuisance training and evaluation by independent branch can protect
one dependence boundary; it cannot remove the shared calendar shock. The
sampling and covariance argument must follow the intended deployment and
observation law.

\section{What assistance and software reuse can responsibly achieve}
\label{exec14:l5-handoffs}

These three pathways suggest distinct roles for an assistant. In sequential
servicing, it can translate a verbal rule into a decision-time history, locate
post-action variables, and ask why assignment and retention are conditionally
ignorable. In a branch rollout, it can distinguish cohort time from event time,
construct candidate comparison groups, and expose anticipation or account
composition changes. In observation construction, it can extract candidate
definitions and mappings while retaining conflicting evidence and unavailable
truth. A generic request to ``run causal inference'' is insufficient for all
three. The proposed product should make the pathway and its resulting target
visible to the analyst before selecting an estimator.

The point-in-time OCI workflow demonstrates a narrower literature precedent:
an accountable human principal, actor/critic assistance, reproducible notebooks,
cross-fitted AIPW and overlap diagnostics under supplied unconfoundedness
\cite{exec13_data_oci}. It supports discussion of human oversight and
point-treatment workflow. It does not establish sequential exchangeability,
identify a rollout counterfactual, or qualify ownership reconstruction. These
are candidate extensions requiring the corresponding scientific and numerical
arguments. A reviewer should be able to see whether the assistant asks for
the missing premise, changes the target explicitly, or refuses the requested
interpretation. An apparently complete notebook can still be wrong if its
target differs from the business question.

The inspected R package \texttt{ltmle} offers a relevant longitudinal
component. It estimates means specific to treatment and censoring, supports
dynamic treatment rules, and returns outcome-regression and weighting
outputs \citep{sw_ltmle}. Its interface
uses column order to encode model chronology and binary treatment nodes.
An adapter would need to validate the actual event and availability clocks
before constructing that order. A continuous contact dose cannot be passed
as though it were binary without an explicit intervention redesign.
Absorbing event coding is appropriate for a one-time event only when its
semantics are actually absorbing. Automatically truncating net outcomes at a
declared range can alter the quantity being averaged, whereas a known affine
scaling can be reversed for an intervention mean. These interface choices
have scientific consequences, rather than merely formatting consequences.
The saved source inspection motivates restricted reuse; it is not an execution
or qualification of the package in a bank setting.

A matched baseline is an analyst using explicit decision interviews, a graph
editor, and deterministic estimation software with the same sources and
information budget. Assistance should be assessed on preserved targets,
unsupported assumptions, actual tool completion, accurate uncertainty, and
review effort. Comparisons should include unavailable histories, unsupported
actions, contaminated controls, missing frames, and labels without gold truth.
An appropriate unresolved answer may be more valuable than a confidently
generated estimate. Genuine improvement would mean reducing analyst effort
while preserving those decisions and making more errors detectable, rather
than increasing the number of successful-looking reports.

The immediate prototype can use the finite mechanisms above and suitable
public panels to test those transformations economically. Known synthetic
truth tests target preservation and numerical arithmetic; a public observed
panel tests the discipline of stating a counterfactual without supplying its
truth by construction. A later bank experiment must resolve the actual
assignment, observation, operating-cost and acceptance questions. The proposal
can justify that experiment through its scoped comparisons and explicit
failure modes. It cannot replace the experiment's missing evidence with the
existence of an agent architecture.
\label{exec14:l5-development-end}
